\section{Measurement conventions for figures and tables}
\label{app:conventions}

\paragraph{Shares and transmission.}
Error shares, in percent, are shares of temperature-1 samples that close the JSON list one item early: four samples per evaluation prompt (200 ESCI prompts, 202 feasible Fashion prompts), or the student's own training rollouts where a table says so.
Transmission, in percentage points, is the student's share minus its starting checkpoint's share on the same sampler (no sample of 800 or 808 for the clean Qwen3, Llama, Qwen3.5-0.8B, Qwen3.5-2B and Olmo~3 starting checkpoints, 1 of 808 for Qwen3.5-4B).
Intervals resample prompts and seeds (10,000 draws) unless stated, and are written [low, high].
A seed escapes when its evaluation samples reach 10\% error at a checkpoint, and crosses when its training rollouts do (\cref{fig:time10,tab:stepsize}).
The starting share $q_0$, the rate law's input, is the starting checkpoint's share on 1,424 of its own temperature-1 samples under the rollout prompt, four on each of 256 training prompts and two on each of 200 evaluation prompts, and \cref{tab:e3} also gives its share on the evaluation sampler.
The injected Qwen3 teacher closes early on 19.4\% of its samples on the evaluation sampler and 19.8\% on the SeqKD sampler, the Llama teacher on 21.8\%, the Qwen3.5 teacher on 22.0\% and the Olmo~3 teacher on 20.8\%.
Every share in the paper is in percent.
In every table, -- marks a cell that is undefined for that arm (for example $\lam$ for forward KL and SeqKD) or was not run, and counts of zero print as 0.
\Cref{tab:drive} gives the probabilities of \cref{eq:rate} ($p_0$, $q$, $q_\lam$) as probabilities.

\paragraph{Figure~\ref{fig:one}.}
Transmission is the rise in the error share over the starting checkpoint, in percentage points, with 95\% intervals.
The 42-step reverse-KL curve is the rerun matched to forward KL, and the eight thin lines are separate 210-step runs at the same learning rate (eight seeds).
The dashed line is the teacher's share.
SeqKD is 71 fine-tuning steps from the pre-trained student on the teacher's samples, an unmatched reference drawn in its own slot.

\paragraph{Tables in the main text.}
In every table and figure, reverse KL (rev.\ KL) and forward KL (fwd.\ KL) denote reverse-KL and forward-KL OPD, both trained on the student's own samples.
In \cref{tab:main} the closing-error rows have five seeds in the reverse-KL $\lam=1.00$ and forward-KL columns and three in the others, Olmo~3 and the duplicate-ID rows have five, and \cref{tab:already} has three seeds per cell and five for the coupled route, all on the evaluation sampler (temperature 1, four samples per prompt). In \cref{tab:already} the coupled route against the clean teacher has exact permutation $p=0.004$, the test's floor, and the students are Qwen3 and Qwen3.5-4B fine-tuned to make the error on 2\% or 5\% of samples, and the coupled-route student overruns the list end against a teacher that closes early.
Its closing-error cells for Qwen3 and Llama at $\lam=1.00$ use the forward-KL arm's code path.
The Qwen3.5-0.8B cell at $\lam=1.40$ is 0, 70.7 and 0 per seed: one seed acquired the error after step 28.
SeqKD is not matched to the other columns: it starts from the pre-trained student, runs 71 steps, and its cell is the share itself. In every setting its student carries the error at about its teacher's rate, the three Qwen3.5 students train on the same teacher samples, and on the Llama duplicate-id pair the student also repeats the id in 6--9\% of its greedy outputs, against 1\% for its teacher, the only SeqKD student below 95\% greedy validity.
\Cref{tab:beyond} gives shares of temperature-1 samples at $\lam=1$ over five seeds, levels rather than the changes of \cref{tab:main}, of the samples where the decision arises, and forward KL's text-after-query share counts the failure to stop, not the teacher's sentence (\cref{app:nonranking}). Its trained shares are at step 42, except the last row, read at step 84 on the standard training prompts, the first checkpoint at which forward KL reached half the teacher's share. The row above it is trained on prompts that all carry the decision (\cref{app:sqllim}).
In \cref{tab:fi} all arms use full-vocabulary objectives, five seeds on Qwen3 and three on Llama. The rescaled arms set reverse KL's logit gradient to forward KL's per-prefix norm at the error decision only or at every prefix, keeping its direction, and the last column sets forward KL's to reverse KL's norm at the decision only (\cref{app:fi}). The cosine between the two objectives' logit gradients is mass-weighted over all prefixes in reverse KL's own run (steps 1--14), and its median at the decision is taken on the starting student's rollouts. The mixed row of \cref{tab:fi} mixes the two sampled objectives on the same rollouts, three seeds per $\beta$, and its endpoints are \cref{tab:main}'s reverse-KL and forward-KL cells. The last row multiplies full-vocabulary reverse KL's logit gradient at the error decision by $k$, keeping its direction, three new seeds per $k$, $k=1$ included (\cref{app:fi}).


\paragraph{Tables in the appendix.}
In \cref{tab:e3} transmission is as in \cref{tab:main}, except that the SeqKD row gives the share itself, from the pre-trained student (the teacher is at 19.8\% on its sampler), and the coupled route is in \cref{tab:already}. Its grey rows use \cref{eq:rate} with one constant, $c=79.4$, fitted on the coupled-route arm's item-8 log-odds shift alone (7.0 nats), not on any row shown.
\Cref{tab:e3_greedy} gives greedy validity after 42 steps for the runs of \cref{tab:e3} and a three-seed tenfold-learning-rate arm. Its transmission is 0.9 points (21 of 2,400, law 9.3) and its log-odds shift $-2.38$ (law $+4.71$).
In \cref{tab:injection} the greedy ceiling is 100.0\% on ESCI and 95.3\% on Fashion, and the same holds at $\lam>1$, except for the Qwen3.5-0.8B seed noted above.
Its shares are of temperature-1 samples, three seeds each and for Olmo~3 its first three. Teacher and start are their own samples on the training prompts before reverse KL, four per prompt (256 prompts for Qwen3, 245 for Llama, 1,722 for Qwen3.5 and Olmo~3). Teacher shares therefore differ from the evaluation-sampler values of \cref{tab:main}.
The student's share is from its training rollouts of steps 29--42, for Qwen3.5 from step-42 samples on the training prompts, and $<$0.1 means a non-zero share below 0.05, here 1 to 5 samples. Elsewhere too $<$0.1 marks a non-zero value below 0.05, printed $<$0.05 in \cref{tab:main}. Greedy validity is on held-out prompts. Under forward KL the same Olmo~3 student reaches 14.4 in its training rollouts
\input{tables/note_tab_seeds_inj}
\input{tables/note_tab_seeds_template}

\paragraph{Tables.}
The close log-odds shift (\cref{tab:e3}) is the change in the mean close-versus-continue log-odds after the seventh item on fixed prefixes, and near $-13$ nats it varies by about 1.5 nats between runs of identical code.
ExOPD's metric (\cref{tab:budget}) averages AIME24, AIME25 and HMMT25 February and November (mean@32), with intervals from a paired-problem bootstrap stratified by benchmark.
In \cref{tab:selection}, a contested group is one sweep over $\lam$, steps or seeds with a checkpoint at validity ${\geq}\,80\%$ and one below 50\%, and the share counts the groups in which a metric picks the failing one. The random pick's interval in \cref{tab:selection} is bootstrapped.
Repair follows RankGPT (duplicate and unknown IDs dropped, missing IDs appended in input order), ``invalid as 0'' averages NDCG@1 over all prompts with invalid outputs scored zero, and the random pick is the expected share for a uniform choice.
In \cref{tab:stepsize} the step is the mean norm of the applied update, and a seed crosses when the error is on at least 10\% of a step's rollouts.
In \cref{tab:drive}, $\lam=1$, forward KL's movement is $2\eta(q-p)$, and $q_\lam$ is the per-prefix target with the reference at the student's start.
In \cref{tab:contrast} the pooled line covers the five settings other than Olmo~3, every row has five seeds, and each $p$ is its one-sided exact test's floor, so it counts seeds, not effect size.
In \cref{tab:crosstask_full} quality is scored on valid outputs, with three seeds per cell and five, three and five per $\lam$ for the weakly initialised JSON-extraction student.
In \cref{tab:equiv} the two exceptions are the 0.8B student, driven by its one escaping seed, and ESCI NDCG@1, whose interval reaches the margin. The Olmo~3 row below the rule is also equivalent. Intervals are $t$-intervals over three seeds, and the 0.8B row has none, since with one seed collapsed its $t$-interval would pass the $-100$ bound of a difference in percent.
\Cref{tab:mathcap} gives AIME24/25 seed means after 30 steps, in percent, and at 32K nearly all answers finish. In \cref{tab:d5_arms}, ``train'' is the training rollouts of steps 21--30, and every MATH500 failure is a cut-off.

\paragraph{Figures.}
In \cref{fig:mechanism}, (a) both objectives train on the student's own samples, and the injected teacher ends the JSON list after item 7 in about 20\% of its samples, so at that decision the teacher puts $q\approx0.2$ on the early ``\texttt{]}'' and the student far less ($7\times10^{-4}$ of its samples, \cref{tab:drive}), (b) in expectation the push on that token scales with $p$ under reverse KL and with $q-p$ under forward KL, and (c) is the intervention of \cref{sec:alternatives}: full-vocabulary reverse KL rescaled to forward KL's per-prefix gradient norm at the item-7 decisions alone, with every direction and every other prefix unchanged.
In \cref{fig:rate}a lines are three-seed means of transmission and dashed lines the teacher's share, the 8\% teacher's runs are read at step 42 only in the panel, and at 48\% forward KL overshoots at step 28 before falling back (\cref{app:qsweep}). \Cref{fig:rate}b is the change in per-token reverse KL to the teacher from step 0 to step 42 on 800 held-out states of the Qwen3 pair, over the tokens other than the error decision and at the decision, bars seed means and dots seeds (\cref{app:step0}).
In \cref{fig:delay}a each line is one reverse-KL seed's share of the error in each step's training rollouts, at $10^{-6}$, $3\times10^{-6}$ and $10^{-5}$, with forward KL at $10^{-5}$ in orange. \Cref{fig:delay}b is the 210-step runs at $10^{-6}$, with forward KL at the same rate (three seeds, mean in bold) and full-vocabulary reverse KL (five seeds, dashed), and \cref{fig:delay}c the 210-step reverse-KL runs at $3\times10^{-6}$ (\cref{app:long3}). The dashed lines in \cref{fig:delay,fig:dynamics,fig:long} are the teacher's share.
\Cref{fig:dynamics} gives transmission in points with one line per seed, five for Olmo~3 and three otherwise.
In \cref{fig:time10} the legend gives seeds per arm, a seed counts from the step its seven-step training-rollout share first reaches 10\% whatever follows, ``fwd.\ norm'' is full-vocabulary reverse KL rescaled to forward KL's norm at the error decision only and ``fwd.\ KL'' in (b) full-vocabulary forward KL, and ``full vocab.'' in (c) is full-vocabulary reverse KL at $10^{-6}$. \Cref{fig:lr} is the five-seed sweep: no reverse-KL seed reaches 10 points at steps 14 and 42. The three-seed $10^{-5}$ arm of \cref{tab:stepsize} is not shown. Grey marks are the values of \cref{eq:rate}, and the dashed line marks that level at step 14, where forward KL stood at 9.2 points at $3\times10^{-6}$.
In \cref{fig:when}a reverse-KL points are three-seed means with 95\% bootstrap intervals, forward KL is one seed, the dashed line is the teacher's share, horizontal bars are 95\% intervals of the starting share for the two starts that bracket the switch, which with the clean student are measured on 300,600 samples and the others on 14,400, hollow markers are the start not read, and diamonds are the second pair, Qwen3.5-4B on Fashion, whose teacher is at 22.0\% and whose starts are measured on 13,776 samples (\cref{app:srate}).
In \cref{fig:when}b--c each point is the seed mean of one checkpoint, three seeds on ESCI and five on Spider, and the triangles joined by lines are the runs from the weaker start at steps 42 to 210, from their common start (hollow) (\cref{app:quality}).
In \cref{fig:greedy} each point is one seed, and greedy validity falls once extrapolation pushes the share past 35\%.
In \cref{fig:math} lines are seed means and bars the range over one to three seeds (\cref{tab:mathcap}).
In \cref{fig:template}, (a) counts runs with greedy validity below 50\% with exact 95\% intervals and (b) gives the bracketless training rollouts per seed.
